\section{Introduction}
\label{sec:intro}

Video games are a compelling testbed for large language model (LLM) and vision-language model (VLM) agents~\cite{berner2019dota, xu2025agents} because they demand the same core capabilities as general-purpose agents~\cite{wang2023voyager, raad2024scaling}: games are \textbf{interactive} and \textbf{long-horizon}, requiring the agent to \textbf{deploy skills over many turns} and to engage in \textbf{reasoning-heavy} decision-making---summarizing the state, inferring intentions, selecting relevant skills, and choosing actions---in scalable, controllable environments, with broader implications for continual-learning agents in real-world applications~\cite{hu2024survey}. At the same time, learning from interaction in such settings is difficult~\cite{liao2025think, lu2025cultivating}: rewards are delayed~\cite{liu2025agentic}, strategies are compositional, and high-quality human demonstrations or skill annotations are scarce. As a result, many existing game-playing agents still rely on high-quality annotated human data~\cite{hu2023language}, or require long training times with human or external-LLM feedback.

Recent progress in self-improving LLM and VLM agents~\cite{li2025self, wang2025vision, he2025visplay} suggests a path beyond annotation-heavy pipelines: agents can learn from large amounts of unlabeled interaction data and reduce reliance on high-quality demonstrations~\cite{gao2025survey, xiangsystematic}. However, long-horizon settings remain difficult due to sparse, delayed rewards, which make policy improvement hard to track and stabilize. Since vanilla self-improvement is often insufficient, memory mechanisms for experience reuse and sustained improvement become essential~\cite{xu2025mem, sarch2024vlm}: the agent must organize, retain, and reuse useful experience across episodes rather than relying only on parametric updates. Recent work further advances this direction by letting agents manage memory and \emph{skills}---reusable, temporally extended behaviors for recurring subgoals---showing that structured memory operations and active updating can improve performance on complex long-horizon tasks in web and embodied agents~\cite{xia2026skillrl, zhang2026memrl, zhang2026memskill}.

Compared with many agent settings, video game episodes are often longer-horizon and more dynamic, with partial observability and uncertainty adding difficulty~\cite{berner2019dota, raad2024scaling}. A single episode may require the agent to deploy \emph{multiple} skills at different stages, each spanning many turns. Treating an episode as a flat trajectory therefore hinders experience reuse: useful behaviors are temporally entangled, context-dependent, and hard to retrieve or improve in isolation. This makes \textbf{skill segmentation} essential---the agent must decompose long trajectories into reusable skill-level units for memory, retrieval, and continual improvement. Despite its importance, skill segmentation for LLM/VLM agents in long-horizon gameplay remains underexplored; changing game states, branching outcomes, and sparse feedback make it substantially harder than in shorter or more structured settings.

In this paper, we propose a \textbf{multi-agent co-evolution} framework for long-horizon game playing that closes the loop between (i) a vision-language \textbf{decision agent} that plays games using primitive actions and tool calls (\eg, \texttt{QUERY\_SKILL}, \texttt{CALL\_SKILL}, \texttt{QUERY\_MEM}) and (ii) an \textbf{agent-managed skill bank} that performs unsupervised skill discovery, maintenance, and retrieval. The decision agent consumes the skill bank for retrieval-augmented planning (protocols) and reward shaping (effect contracts); the skill bank agent ingests the decision agent's trajectories and turns them into segmented, labeled skills with protocols and contracts via a multi-stage pipeline (boundary proposal, segment decode, contract learning, bank maintenance). Both agents are trained in the loop: the decision agent with GRPO and the skill bank with Hard-EM and LoRA adapters, so that better rollouts yield better skills and a better skill bank yields better decision-making. To the best of our knowledge, this is the first framework to jointly couple VLM-based game decision-making with an agentic skill-bank pipeline for unsupervised long-horizon skill segmentation and continual skill refinement in a unified co-evolution loop.

Our main contributions include:
\begin{itemize}
    \item We present a general co-evolution framework that closes the loop between gameplay and skill acquisition: a VLM decision agent interacts with the environment to collect trajectories, and an agent-managed skill pipeline converts them into reusable skills (protocols and effect contracts) that improve future gameplay via retrieval and reward shaping.
    \item We propose a practical agentic pipeline for skill bank maintenance that transforms unlabeled trajectories into reusable skills via segmentation, contract learning, and bank maintenance. The skill bank is continually updated, allowing the decision agent to retrieve relevant skills and memories via tool calling for retrieval-augmented decision-making across diverse game environments.
    \item We conduct extensive experiments across multiple long-horizon video game environments, showing that the proposed co-evolution framework outperforms vanilla VLM agents, non-co-evolving methods, and training-free baselines, with a notable \wxy{XX}\% improvement, validating the contribution of each component to long-horizon performance.
\end{itemize}
